\section*{第 \ref{ch:g7:prop} 章——比例}

\begin{solution}{exo:g7:prop:1}
第一张表：商 $\frac{7.5}{3} = 2.5$，
$\frac{17.5}{7} = 2.5$，$\frac{27.5}{11} = 2.5$：\omterm{def:g6:propdata:def}{成正比例}，
比例系数为 $2.5$。

第二张表：$\frac62 = 3$ 与 $\frac{15}{5} = 3$，但
$\frac{28}{9} \approx 3.1$：不\omterm{def:g6:propdata:def}{成正比例}。
\end{solution}

\begin{solution}{exo:g7:prop:2}
用交叉相乘法：$5 \times p = 6 \times 8$，故
$p = \frac{48}{5} = 9.60$ 欧元。
\end{solution}

\begin{solution}{exo:g7:prop:3}
速率：$36 \div 3 = 12$ 页/分钟。$5$ 分钟内：
$60$ 页。对于 $96$ 页：$96 \div 12 = 8$ 分钟。
\end{solution}

\begin{solution}{exo:g7:prop:4}
$30\,\%$ 的 $250$：$0.3 \times 250 = 75$。\quad
$15\,\%$ 的 $60$：$0.15 \times 60 = 9$。\quad
$t\,\%$ 的 $80$ 是 $20$：$\frac{20}{80} = \frac{25}{100}$，故
$t = 25$。
\end{solution}

\begin{solution}{exo:g7:prop:5}
$\frac{18}{25} = \frac{18 \times 4}{25 \times 4} = \frac{72}{100} =
72\,\%$.
\end{solution}

\begin{solution}{exo:g7:prop:6}
路线：$9 \times 25\,000 = 225\,000$ cm $= 2.25$ km。
湖：实地距离 $2$ km $= 200\,000$ cm，故
图上距离为 $200\,000 \div 25\,000 = 8$ cm。
\end{solution}

\begin{solution}{exo:g7:prop:7}
$4.3$ m $= 430$ cm，而 $430 \div 43 = 10$ cm。
\end{solution}

\begin{solution}{exo:g7:prop:8}
每份巧克力：$240 \div 8 = 30$ g。$300$ g 时：
$300 \div 30 = 10$ 份，恰好\omterm{def:g6:wholes:divisible}{整除}。
\end{solution}

\begin{solution}{exo:g7:prop:9}
\emph{1.} 折扣为 $20$ 欧元，即 $\frac{20}{80} = 25\,\%$ 的
原价。

\emph{2.} 涨价 $20$ 欧元，现与\emph{新的}标准量 $60$ 相比，
$\frac{20}{60} = \frac13 \approx 33\,\%$。
百\omterm{def:g6:fractions:def}{分数}总是指相对于某个整体；整体变了，
百\omterm{def:g6:fractions:def}{分数}也就变了。
\end{solution}

\begin{solution}{exo:g7:prop:10}
由完整的一列得系数：$\frac{17.5}{10} = 1.75$。于是
$y = 4 \times 1.75 = 7$，$x = \frac{10.5}{1.75} = 6$。
\end{solution}

\begin{solution}{exo:g7:prop:11}
\emph{1.} 复印一次后：$10 \times 0.8 = 8$ cm。复印两次后：
$8 \times 0.8 = 6.4$ cm。

\emph{2.} 第二次缩小是作用于已缩小的复印件，
而不是原图：两个系数相乘，
$0.8 \times 0.8 = 0.64$，故两次缩小后为
原图的 $64\,\%$，而非 $60\,\%$。
\end{solution}

\begin{solution}{pb:g7:prop:1}
\textbf{1.} 高度与影长\omterm{def:g6:propdata:def}{成正比例}，故由
交叉相乘法（见\cref{thm:g7:prop:cross}），木棍的一列为
$(1,\ 0.4)$，树的一列为 $(h,\ 3.2)$：
$0.4 \times h = 1 \times 3.2$，因此
$h = \frac{3.2}{0.4} = 8$ m。

\textbf{2.} 塔高：$\frac{26}{0.4} = 65$ m。孩子的影长：
$1.5 \times 0.4 = 0.6$ m（影长 $=$ 高度 $\times$ 此刻的
比例系数 $0.4$）。

\textbf{3.} 太阳在天空移动时，
高度与影长之间的比例系数随之改变；只有同一时刻测得的数据
才具有相同的比例系数。

\textbf{4.} 此时比例系数为 $1$：每个高度
都等于其影长。因此金字塔\omterm{def:g5:solids:def}{顶点}
高 $147$ m——即大金字塔的高度（只需
测出影子末端，自\omterm{def:g5:solids:def}{顶点}正下方的点量起即可）。

\textbf{5.} 列的商
$\frac{\text{影长}}{\text{高度}}$——即
此刻的比例系数——对每个竖直
物体都相同：这正是使身高与影长的表格
成为\omterm{def:g7:prop:table}{比例表}（见\cref{def:g7:prop:table}）。

\textbf{6.} 太阳光线\omterm{def:g6:lines:perp}{平行}射来。若地球是\emph{平的}，
两条\omterm{def:g6:lines:perp}{平行}光线将以相同角度照射两根竖直的木棍，
两者的影子就\omterm{def:g6:propdata:def}{成正比例}：要么都无影，要么都有影。
一根木棍无影（塞恩）、一根有影（亚历山大），
在同一时刻出现这种情况，
在平坦的地球上是不可能的：两根竖直线
必须指向不同方向——地表是弯曲的。

\textbf{7.} 图中，塞恩处的竖直线笔直指向太阳，
亚历山大处的竖直线则倾斜了
从地心看去两城市方向之间的夹角。由于光线
\omterm{def:g6:lines:perp}{平行}，这一倾斜角恰好就是在亚历山大测得的
光线与木棍之间的夹角 $7.2^\circ$。

\textbf{8.} $\frac{7.2}{360} = \frac{72}{3600} = \frac{1}{50}$：
即整圈的五十分之一。

\textbf{9.} 弧长与角度成正比：若
$7.2^\circ$ ——即整圈的五十分之一——对应
$800$ km，则整圈对应
\[
50 \times 800 = 40\,000 \text{ km}.
\]

\textbf{10.} \omterm{def:g4:geometry:circle}{直径} $= \frac{\text{周长}}{\pi}
\approx \frac{40\,000}{3.14} \approx 12\,700$ km ——与
现代值 $12\,742$ km 相比：误差远小于百分之一，
仅用一根木棍即可。

\textbf{11.} $40\,000$ km $= 4\,000\,000\,000$ cm，而
$4\,000\,000\,000 \div 40\,000\,000 = 100$ cm $= 1$ m 的
\omterm{def:g6:measure:perimeter}{周长}。\omterm{def:g4:geometry:circle}{直径}：$100 \div 3.14 \approx 32$ cm ——
一个漂亮的桌面地球仪。

\textbf{12.} $4.8$ km $= 480\,000$ cm，而
$480\,000 \div 40\,000\,000 = 0.012$ cm $= 0.12$ mm。阿尔卑斯山的最高峰，
在一个\omterm{def:g6:measure:perimeter}{周长}一米的地球仪上，仅相当于十分之一毫米，
不过一粒尘埃。以此\omterm{def:g7:prop:scale}{比例尺}看，地球
比大多数打磨过的球还要光滑。

\textbf{13.} $40\,000 \div 40 = 1\,000$ 天，而
$1\,000 \div 365 \approx 2.7$：几乎要走上三年。

\textbf{14.} $\frac{10}{6\,370} \approx 0.0016$，即约
$0.2\,\%$（更精确地说是 $0.16\,\%$）：大气层
不过是覆在星球表面的一层薄薄的皮。

\textbf{15.} 误差：$42\,000 - 40\,000 = 2\,000$ km，而
$\frac{2\,000}{40\,000} = \frac{5}{100} = 5\,\%$。一句话概括：
仅凭一根木棍、一口井、一条测量过的道路和一个比例关系，
埃拉托斯特尼就把一颗星球的大小测到了几个百分点以内——
数学用很少的东西，便能走得很远。
\end{solution}
